\documentclass[12pt]{article}

%% Basic document formatting
\usepackage{amsmath, amsthm, amssymb, amsfonts}
\usepackage{mathtools}
\usepackage{xspace}
\usepackage{thmtools}
\usepackage{graphicx}
\usepackage{setspace}
\usepackage{fancyhdr}
\usepackage{titling}
\usepackage[left=0.4in,right=0.4in,top=1in,bottom=1in]{geometry}
\usepackage{float}
\usepackage{hyperref}
\usepackage{mathrsfs}
\usepackage{tabularx}
\usepackage[utf8]{inputenc}
\usepackage[english]{babel}
\usepackage{framed}
\usepackage[dvipsnames]{xcolor}
\usepackage{environ}
\usepackage{tcolorbox}
\tcbuselibrary{theorems,skins,breakable}

\usepackage{enumitem}
\setlist[enumerate]{leftmargin=*}
\setlist[enumerate,1]{labelindent=\parindent}
\setlist[enumerate,2]{labelindent=0pt}

% Blackboard Bold
\newcommand{\Z}{\mathbb{Z}}                 % integers
\newcommand{\Zpl}{\mathbb{Z}_{+}}           % positive integers
\newcommand{\Znn}{\mathbb{Z}_{\geq 0}}      % nonnegative integers. 
\newcommand{\Q}{\mathbb{Q}}                 % rationals
\newcommand{\Qpl}{\mathbb{Q}_{+}}           % positive Rationals
\newcommand{\R}{\mathbb{R}}                 % reals
\newcommand{\Rpl}{\mathbb{R}_{+}}           % positive Reals
\newcommand{\C}{\mathbb{C}}                 % complex numbers
\newcommand{\F}{\mathbb{F}}                 % field

% Words
\newcommand{\st}{\text{ such that }}
\newcommand{\wrt}{\text{ with respect to }}
\newcommand{\with}{\text{ with }}
\newcommand{\ie}{\text{, i.e. }}
\newcommand*{\ditto}{\text{---}\texttt{"}\text{---}}

% Operators
\newcommand{\abs}[1]{\left|#1\right|}                   % absolute value
\newcommand{\norm}[1]{\abs{\abs{#1}}}
\newcommand{\floor}[1]{\left\lfloor #1 \right\rfloor}   % floor
\newcommand{\ceil}[1]{\left\lceil #1 \right\rceil}      % ceiling
\newcommand{\powset}[1]{\mathscr{P}(#1)}

% Category Theory 
\renewcommand{\hom}{\text{Hom}}
\newcommand{\homspace}[3]{\hom_{#1}(#2, #3)}
\newcommand{\coproduct}[9]{
  \begin{tikzcd}[ampersand replacement=\&]
      \& \& #1 \arrow[dll, bend right, "#6"] \arrow[dl, "#5"] \\
   #4 \& #3 \arrow[l, "#9"] \\
      \& \& #2 \arrow[ull, bend left, "#8"] \arrow[ul, "#7"]
  \end{tikzcd}
}

\newcommand{\product}[9]{ 
  \begin{tikzcd}[ampersand replacement=\&]
      \& \& #3 \\
      #1 \arrow[r, "#7"] \arrow[urr, bend left, "#5"] \arrow[drr, bend right, "#6"] \& #2 \arrow[ur, "#8"] \arrow[dr, "#9"] \\ 
      \& \& #4
  \end{tikzcd}
}


% Algebra 
\newcommand{\Syl}[2]{\text{Syl}_{#1}{(#2)}}           % set of Sylow-p subgroups 
\newcommand{\<}{\langle}                % \<x\>, subgroup generated by x 
\renewcommand{\>}{\rangle}
\renewcommand{\index}[2]{[#1 : #2]}
\newcommand{\id}{\text{id}}             % identity element
\newcommand{\lhdneq}{%
  \mathrel{\ooalign{$\lneq$\cr\raise.22ex\hbox{$\lhd$}\cr}}}
\newcommand{\order}[1]{\text{o}(#1)}    % order of an element
\newcommand{\spec}{\operatorname{Spec}}
\newcommand{\rad}{\operatorname{rad}}   % radical of an ideal 
\newcommand{\sym}[1]{\mathfrak{S}_{#1}}
\newcommand{\aut}[1]{\text{Aut}(#1)} 
\newcommand{\gal}[2]{\text{Gal}(#1 / #2)} 
\newcommand{\inn}[1]{\text{Inn}(#1)} 
\let\oldcong\cong
\let\oldequiv\equiv
\renewcommand{\cong}{\oldequiv}
\renewcommand{\equiv}{\oldcong}

\newcommand{\triangleleftneq}{\mathrel{\ooalign{%
  $\trianglelefteq$\cr
  \hidewidth\raise0.06ex\hbox{$\not\mkern4mu$}\hidewidth\cr
}}}

\makeatletter % cycle
\newcommand{\cyc}[1]{(\mathbf{\cyc@process#1\relax})}
\def\cyc@process#1#2\relax{%
	#1%
	\ifx\relax#2\relax
	\else
		\,\cyc@process#2\relax
	\fi
}
\makeatother

\newcommand{\GL}[2]{\text{GL}_{#1}(#2)}                             % general linear group
\newcommand{\SL}[2]{\text{SL}_{#1}(#2)}                             % special linear group
\newcommand{\gl}[2]{\mathfrak{gl}_{#1}(#2)}
\renewcommand{\sl}[2]{\mathfrak{sl}_{#1}(#2)}
\newcommand{\so}[2]{\mathfrak{so}_{#1}(#2)}
\newcommand{\su}[1]{\mathfrak{su}(#1)}

% Linear Algebra
\newcommand{\bmat}[1]{\begin{bmatrix}#1\end{bmatrix}}   % bracketed matrix
\newcommand{\rank}{\operatorname{rank}}                 % rank
\newcommand{\nullity}{\operatorname{nullity}}           % nullity
\newcommand{\tr}{\operatorname{tr}}

% Combinatorics
\newcommand{\qnum}[1]{\left[#1\right]_q}                % q-analog
\newcommand{\qfact}[1]{\left[#1\right]!_q}              % q-analog factorial 
\newcommand{\qbinom}[2]{\binom{#1}{#2}_q}               % Gaussian binomial coefficient

% Topology/Analysis
\newcommand{\ball}[2]{\text{B}_{#1}(#2)}  % B_r(x): open r-balls around x
\newcommand{\diam}{\text{diam}}           % diamter of a set in metric space 
\newcommand{\lowint}[2]{\underline{\int_{#1}^{#2}}}
\newcommand{\upint}[2]{\overline{\int}_{#1}^{#2}}

% Misc Notation
\newcommand{\oneton}[1]{\{1, \ldots, #1\}}
\newcommand{\defeq}{\vcentcolon=}   % :=
\newcommand{\eqdef}{=\vcentcolon}   % =: 
\renewcommand{\bf}[1]{\textbf{#1}}
\newcommand{\im}{\operatorname{im}}
\newcommand{\fnrest}[2]{\left.#1\right|_{#2}}
\newcommand{\isom}{\overset{\sim}{\rightarrow}} 


\renewcommand{\thesection}{\arabic{section}.}
\renewcommand{\thesubsection}{(\alph{subsection})}

\newtheorem{claim}{Claim}
\newtheorem{theorem}{Theorem}
\newtheorem{proposition}{Proposition}
\newtheorem*{lemma}{Lemma}

%% Headers & title setup
\newcommand{\course}{Math140C}
\newcommand{\myname}{Jay Ser}
\setlength{\headheight}{14.5pt}
\pagestyle{fancy}
\fancyhf{}
\renewcommand{\headrulewidth}{0.4pt}
\lhead{\course}
\rhead{\myname}
\cfoot{\thepage}
\setlength{\droptitle}{-4em} 
\title{\course\ - HW \#1}
\author{\myname}
\date{2026.04.12} 

\begin{document}
\maketitle
\thispagestyle{fancy}

%------------------------------------------------------------------------------%

\section*{Q5}
Let $A = \bmat{a_1 & \ldots & a_n}$. 
For any $x = \bmat{x_1 \\ \vdots \\ x_n} \in \R^n$, 
\begin{equation} \label{eq:5_A_rep}
  Ax = \sum_{i = 1}^{n} {a_i x_i}.
\end{equation}
So just let $y = \bmat{a_1 \\ \vdots \\ a_n}$, then $Ax = x \cdot y$ for all $x \in \R^n$. 
It's clear this particular $y$ is the only coordinate vector such that the condition $Ax = x \cdot y$ is satisfied because \eqref{eq:5_A_rep} shows that each coordinate of $A$ contributes to the value $Ax$. 

It has already been shown in class that $\abs{Ax} \leq \abs{A} \abs{x}$ for any $x \in \R^n$. 
Notice that $\abs{A} = \abs{y}$. 
So it is clear that $\norm{A} \leq \abs{y}$. 
To prove the reverse inequality, consider the normalized vector $x = y / \abs{y}$. 
Then $\abs{x} = 1$ and 
$$\abs{Ax} = \abs{\frac{1}{\abs{y}} \sum a_{i}^{2}} = \frac{1}{\abs{y}}\abs{y}^2 = \abs{y},$$
hence $\abs{y} \leq \norm{A}$.  

\section*{Q6} 
If $x \neq 0$ and $y \neq 0$, then one can just calculate the partial derivatives as one does in elementary calculus without doing any weird limit things: 
\begin{align*}
  (D_1 f)(x, y) & = \frac{y(y^2 - x^2)}{(x^2 + y^2)^2} \\ 
  (D_2 f)(x, y) & = \frac{x(x^2 - y^2)}{(x^2 + y^2)^2}
\end{align*}
Suppose $y = 0$ and $x \neq 0$.  
Then 
\begin{align*}
  (D_1 f)(x, 0) & = \lim_{t \rightarrow 0} \frac{\frac{(x + t)0}{(x + t)^2 + 0^2} - \frac{x0}{x^2 + 0^2}}{t} = 0 \\ 
  (D_2 f)(x, 0) & = \lim_{t \rightarrow 0} \frac{\frac{xt}{x^2 + t^2}}{t} = \frac{1}{x}
\end{align*}
It is clear by symmetry that the partials exist when $x = 0$ and $y \neq 0$. 

Finally, suppose $x = 0$ and $y = 0$. 
$$(D_1 f)(0, 0) = \lim_{t \rightarrow 0} \frac{0 - 0}{t} = 0$$
and $(D_2 f)(0, 0) = 0$ by the same calculation. 
So the partials of $f$ exist everywhere. 
However, $f$ is not differentiable at $(0, 0)$. 
If $f$ was differentiable at $(0, 0)$, then the matrix $f'$ would be given by the partials at $(0, 0)$, namely 
$$f'(0, 0) = \bmat{0 & 0}.$$
But $A = \bmat{0 & 0}$ does not satisfy the definition of the derivative:
writing $h = (h_1, h_2)$, 
\begin{align*}
  \lim_{h \rightarrow 0} \frac{\abs{f(h) - f(0) - Ah}}{\abs{h}} 
  & = \lim_{h \rightarrow 0} \frac{\abs{h_1}\abs{h_2}}{(h_{1}^{2} + h_{2}^{2})\abs{h}} \\ 
  & = \infty \neq 0,
\end{align*}
where the limit to infinity makes sense since the numerator tends to 0 much faster than the denominator. 

\section*{Q8}
Since $x \in E$ is a local maximum of $f$, $\exists \delta > 0$ such that for all $t \in \R^n$ with $\abs{t} < \delta$, $f(x + h) \leq f(x)$. 
Fix any unit vector $u \in \R^n$ and define $g(t) \defeq f(x + tu)$ for $t \in \R$. 
For all $\abs{t} < \delta$, $\abs{x + tu} \leq f(x)$, so $g$ has a local maximum at $t = 0$. 
$f$ is differentiable at $x$, so $g$ is differentiable at $t = 0$ because $g'(0)$ is the directional derivative of $f$ along $u$. 
Further, since $t = 0$ is a local maximum, 
\begin{align*}
  f'(x)u & = \lim_{t \rightarrow 0} \frac{f(x + tu) - f(x)}{t} \\ 
         & = g'(0) = 0
\end{align*}
Because $f'(x)u = 0$ holds for any unit vector in $\R^n$, $f'(x)$ maps any vector of a normalized basis to 0, which means $f'(x)$ is just the zero map. 

\section*{Q15}
\subsection{}
$$(x^4 + y^2)^2 - 4x^4 y^2 = x^8 + 2x^4 y^2 + y^4 - 4x^4 y^2 = (x^4 - y^2)^2.$$
Obviously the right side is nonnegative, so $(x^4 + y^2)^2 \geq 4x^4 y^2$. 

It's clear that $f$ is continuous away from $(0, 0)$. 
To show that $f$ is continuous at $(0, 0)$ as well, it must be shown that $f(x, y) \rightarrow 0$ as $(x, y) \rightarrow (0, 0)$. 
This is indeed the case: 
by the inequality, 
$\frac{4x^6 y^2}{(x^4 + y^2)^2} \leq x^2$ and $x^2 \rightarrow 0$ as $(x, y) \rightarrow (0, 0)$, so $\frac{4x^6 y^2}{(x^4 + y^2)^2} \rightarrow 0$ too. 
It's clear that the other summands of $f$ tend to $0$ as $(x, y) \rightarrow (0, 0)$. 
 
\subsection{}
\begin{align*}
  g_{\theta}(t) & = t^2 \cos^{2}{\theta} + t^2 \sin^{2}{\theta} - 2t^2 \cos^{2}{\theta} (t\sin{\theta}) - \frac{4t^{6} \cos^{6}{\theta} (t^2 \sin^{2}{\theta})}{(t^4 \cos^{4}\theta + t^2 \sin^{2}{\theta})^2} \\ 
                & = t^2 - 2t^3 \cos^{2}{\theta} \sin{\theta} - \frac{4t^4 \cos^{6}{\theta}\sin^{2}{\theta}}{(t^2 \cos^{4}{\theta} + \sin^{2}{\theta})^{2}} \\ 
\end{align*}
One can see from the quotient rule that the first and second derivatives of the third term have some power of $t$ on the numerator the denominator is nonzero at $t = 0$, so 
$g_{\theta}(0) = g'_{\theta}(0) = 0$ and $g''_{\theta}(0) = 2$. 
Hence $g_{\theta}$ has a strict local minimum at $t = 0$. 

\subsection{}
\begin{align*}
  f(x, x^2) & = x^2 + x^4 - 2x^2 x^2 - \frac{4x^6 x^4}{(x^4 + x^4)^2} \\ 
            & = x^2 - x^4 - \frac{4x^{10}}{4x^8} \\ 
            & = -x^4
\end{align*}
So in any neighborhood around $(0, 0)$, there is some $(t, t^2) \neq (0, 0)$ in that neighborhood such that $f(t, t^2) = -t^4 < 0$, such that $(0, 0)$ is not a local minimum for $f$. 

\end{document}
